\chapter{Ergebnisse und Interpretation}
\label{chap:ergebnisse}

In diesem Kapitel werden die Ergebnisse der statistischen Modellierung präsentiert und im Hinblick auf die Forschungsfrage nach den Risikofaktoren für Armut interpretiert. Die Analyse basiert auf dem transformierten Regressionsmodell mit robusten Standardfehlern (HC3), welches die höchste statistische Belastbarkeit aufweist.

\section*{Identifikation zentraler Risikofaktoren}
Die Regressionsanalyse identifiziert zwei primäre strukturelle Risikofaktoren, die einen signifikanten positiven Einfluss auf die Armutsrate haben:

\begin{itemize}
    \item \textbf{Berufsstatus "`Labour"':} Mit einem Koeffizienten von $0,2339$ ($p < 0,001$) stellt die Zugehörigkeit zur Gruppe der unqualifizierten Tagelöhner das größte Armutsrisiko dar. Dies bestätigt Booths deskriptive Beobachtungen auf mathematischer Ebene.
    \item \textbf{Weibliche Haushaltsvorstände:} Sektionen mit einem hohen Anteil an weiblichen Vorständen weisen einen signifikanten Koeffizienten von $0,1583$ ($p = 0,004$) auf. Dies deutet darauf hin, dass Haushalte ohne männlichen Verdiener (oft bedingt durch Arbeitsunfälle oder Verwitwung) im viktorianischen London besonders vulnerabel waren.
\end{itemize}

\section*{Schutzfaktoren und soziale Stabilität}
Im Gegensatz dazu konnten deutliche Schutzfaktoren identifiziert werden, die mit einer signifikanten Verringerung der Armutsrate korrelieren:
\begin{itemize}
    \item Berufe im Sektor \textbf{Salaried} (Angestellte) sowie im Bereich \textbf{Refreshments} (Gastgewerbe) weisen stark negative Koeffizienten auf (z. B. $-0,2185$, $p < 0,001$). Regelmäßige Einkommen und eine feste Anstellung waren somit die effektivsten Mittel zur Armutsprävention.
\end{itemize}

\section*{Demografie vs. Struktur}
Ein wesentliches Finding der Analyse ist die fehlende Signifikanz der demografischen Variablen im multiplen Modell:
\begin{itemize}
    \item \textbf{Kinderzahl:} Entgegen der ursprünglichen Hypothese liefert die Variable \textit{kids\_per\_household} ($p = 0,414$) keinen signifikanten zusätzlichen Erklärungsbeitrag zur Armutsrate, sobald der Berufsstatus kontrolliert wird.
    \item \textbf{Alleinstehende Männer:} Auch das \textit{unmarried\_ratio} ($p = 0,359$) ist statistisch nicht signifikant.
\end{itemize}
Dies deutet darauf hin, dass Armut im 19. Jahrhundert primär ein \textbf{strukturelles Problem des Arbeitsmarktes} und weniger ein rein demografisches Problem der Haushaltsgröße war.

\section*{Güte und Belastbarkeit des Modells}
Das Modell erklärt ca. $23,8\,\%$ der Varianz der Armutsrate ($R^{2} = 0,238$). Während dies für sozialwissenschaftliche Daten ein solider Wert ist, bleibt ein Großteil der Varianz durch Faktoren ungeklärt, die nicht in den historischen Tabellen erfasst wurden, z.b. individuelle Gesundheitszustände oder lokale Mietpreise. Es muss zudem kritisch angemerkt werden, dass es sich um \textbf{Korrelationen}, nicht zwingend um Kausalitäten handelt. 